\section{Problem Statement}

Design and implement a Job Shop Scheduler that assigns operations to machines to minimize the makespan --- the total time from the start of the first operation to the completion of the last. Each job consists of an ordered sequence of operations, where each operation requires a specific machine for a specific duration. A machine can process only one operation at a time, and operations within a job must execute in their prescribed order. Job Shop Scheduling is NP-hard, requiring heuristic or metaheuristic approaches for practical instances.

\section{Core Concepts}

\begin{itemize}
  \item \textbf{Makespan:} The total elapsed time from the start of scheduling to the completion of the last operation across all jobs. This is the primary objective to minimize.
  \item \textbf{Operation Precedence:} Within each job, operations must execute sequentially --- operation $k+1$ cannot start until operation $k$ finishes.
  \item \textbf{Machine Contention:} Each machine can process at most one operation at a time. Operations assigned to the same machine must be serialized.
  \item \textbf{Critical Path:} The longest chain of dependent operations determines the makespan. Identifying and optimizing the critical path is key to reducing overall schedule length.
  \item \textbf{Dispatching Rule:} Heuristic rules such as Shortest Processing Time (SPT), Longest Processing Time (LPT), and First Come First Served (FCFS) for ordering operations on machines.
\end{itemize}

\section{Key Challenges}

\begin{itemize}
  \item \textbf{Solution Representation:} Choosing between operation sequences per machine or explicit start-time assignments, each with different trade-offs for manipulation and evaluation.
  \item \textbf{Feasibility Enforcement:} Ensuring that decoded schedules respect both operation precedence within jobs and single-operation-at-a-time constraints on machines.
  \item \textbf{Makespan Computation:} Efficiently decoding a solution representation into a concrete schedule and computing the resulting makespan.
  \item \textbf{Critical Path Identification:} Determining which operations lie on the critical path and therefore directly influence the makespan.
  \item \textbf{Local Search Moves:} Improving schedules by swapping adjacent operations on the same machine, ensuring swaps maintain feasibility and reduce makespan.
\end{itemize}

\section{Sample Input}

\begin{lstlisting}[style=json, caption={data/input\_sample.json}]
{
  "machines": ["M1", "M2", "M3"],
  "jobs": [
    {
      "id": "J1",
      "operations": [
        {"machine": "M1", "duration": 3},
        {"machine": "M2", "duration": 2},
        {"machine": "M3", "duration": 4}
      ]
    },
    {
      "id": "J2",
      "operations": [
        {"machine": "M2", "duration": 5},
        {"machine": "M3", "duration": 1},
        {"machine": "M1", "duration": 3}
      ]
    },
    {
      "id": "J3",
      "operations": [
        {"machine": "M3", "duration": 2},
        {"machine": "M1", "duration": 4},
        {"machine": "M2", "duration": 3}
      ]
    }
  ]
}
\end{lstlisting}

\vfill
\noindent\rule{\textwidth}{0.4pt}
{\small\textit{For full project requirements, STL restrictions, submission format, and grading criteria, refer to \textbf{base.pdf} (available on the \href{https://github.com/BestSithInEU/cse211-term-projects/releases}{Releases page}).}}
